% Generated draft: requires mathematical and source review.
\documentclass[11pt]{article}
\usepackage[margin=25mm]{geometry}
\usepackage{fontspec}
\setmainfont{Latin Modern Roman}
\usepackage{amsmath,amssymb,mathtools,unicode-math}
\setmathfont{Latin Modern Math}
\usepackage{enumitem,needspace,xurl,hyperref}
\setlength{\parindent}{0pt}
\setlength{\parskip}{5pt}
\setlength{\emergencystretch}{3em}
\newcommand{\N}{\mathbb{N}}
\newcommand{\Z}{\mathbb{Z}}
\newcommand{\Q}{\mathbb{Q}}
\newcommand{\R}{\mathbb{R}}
\newcommand{\C}{\mathbb{C}}
\newcommand{\F}{\mathbb{F}}
\newcommand{\A}{\mathbb{A}}
\newcommand{\PP}{\mathbb{P}}
\newcommand{\E}{\mathbb{E}}
\newcommand{\eps}{\varepsilon}
\newcommand{\dd}{\,\mathrm{d}}
\begin{document}
\section*{UnsolvedMath ID: 700466}
\textbf{Source problem code:} catalogue entry
\subsection{Definitions and mathematical statement}
For a class \(\mathcal K\) of groups, a group \(G\in\mathcal K\) is \emph{existentially closed in \(\mathcal K\)} if every finite system of group equations and inequations with parameters from \(G\) that has a solution in an overgroup \(H\in\mathcal K\) already has a solution in \(G\). A subgroup \(C\leq G\) is characteristic if \(\alpha(C)=C\) for every \(\alpha\in\operatorname{Aut}(G)\). Here the assertion that \(G\) is characteristically simple means that its only characteristic subgroups are \(\{1\}\) and \(G\); the source does not specify whether nontriviality is additionally part of this terminology.

A group is locally nilpotent, locally soluble, or locally finite if every finitely generated subgroup is, respectively, nilpotent, soluble, or finite. It is torsion-free if its identity is its only element of finite order. For a set \(\pi\) of primes, use the conventional meaning of a \(\pi\)-group: every element has finite order and every prime divisor of that order belongs to \(\pi\). With this convention, the classes in part (b) are
\[
\begin{aligned}
L\mathfrak N^{+}&=\{G:G\text{ is locally nilpotent and torsion-free}\},\\
L\mathfrak S_{\pi}&=\{G:G\text{ is locally soluble and a }\pi\text{-group}\},\\
L\mathfrak S^{+}&=\{G:G\text{ is locally soluble and torsion-free}\},\\
L\mathfrak S&=\{G:G\text{ is locally soluble}\}.
\end{aligned}
\]
The source does not specify the allowed choices of \(\pi\), including whether the empty set is excluded; this qualification remains unresolved.

\textbf{Part (b), the remaining questions in the supplied record.} For each of the four classes
\[
\mathcal K\in\{L\mathfrak N^{+},L\mathfrak S_{\pi},L\mathfrak S^{+},L\mathfrak S\},
\]
is the following assertion true?
\[
\forall G\in\mathcal K,\qquad
G\text{ existentially closed in }\mathcal K
\ \Longrightarrow\
\forall C\leq G,\quad
\bigl(C\text{ characteristic in }G\Longrightarrow C=\{1\}\text{ or }C=G\bigr).
\]
Each choice of \(\mathcal K\) gives a separate question, with \(\pi\) held fixed for the second class. There is no countability hypothesis in these questions. The source describes existential closedness as a local property and comments that obtaining global properties from it appears difficult; it also refers to archived part (a).

\textbf{Part (a), supplied archived context.} For a prime \(p\), let \(L\mathfrak N_p\) denote the class of locally finite \(p\)-groups, equivalently groups whose finitely generated subgroups are finite groups of order a power of \(p\). The archived question asks whether every group existentially closed in \(L\mathfrak N_p\) is characteristically simple, with no countability assumption in the question. The supplied text attributes to Berthold Maier, Freiburg, an affirmative answer for countable existentially closed members and the stronger statement that, for fixed \(p\), there is exactly one such countable group up to isomorphism. The supplied archive summary reports a negative answer in general, citing S. R. Thomas (1985). These are supplied historical assertions, not independently verified conclusions of this draft.
\subsection{Short English statement}
For each of four specified classes of locally nilpotent or locally soluble groups, ask whether every existentially closed member has no proper nontrivial characteristic subgroup. Preserve the unspecified restrictions on \(\pi\). Archived part (a) concerns locally finite \(p\)-groups: the supplied record reports an affirmative countable case and a negative general answer. Its overall status is supplied as partially solved, without independent verification.
\subsection{Sources}
\begin{itemize}
\item Supplied local UnsolvedMath registry record: permanent ID \(700466\), exact source problem code \texttt{catalogue entry}, title \emph{Kourovka Notebook Problem 8.14}; upstream revision \texttt{d371ee917ac9647fbb06cd6fc529deec6f3dda82}.
\item O. H. Kegel, proposer, 1982. The supplied source identifies \emph{The Kourovka Notebook (Unsolved Problems in Group Theory)}, 8th issue (1982), Problem 8.14, updated in the 21st issue (2026), arXiv:1401.0300v48. The collection description names E. I. Khukhro and V. D. Mazurov as editors. This is a supplied source identification; the publication has not been independently consulted for this draft.
\item Supplied Archive of Solved Problems entry: negative answer, attributed to S. R. Thomas, \emph{Arch. Math.}, volume 44 (1985), pages 98--109. No article title or external link is supplied.
\item The supplied statement attributes the countable case of part (a), including uniqueness up to isomorphism of the countable existentially closed locally finite \(p\)-group, to Berthold Maier, Freiburg. No separate publication reference is supplied.
\item The supplied status review is dated \texttt{2026-10-16} and labels the record \texttt{partially\_solved}, with solved parts in the archive and remaining parts in the main section. This assessment is supplied dated curation, not independently verified literature assessment.
\end{itemize}
\end{document}
